\section{The Lane Abstraction}

\label{sec:lanes}
%% ============================================================================

The primary innovation of QuasarSTM is replacing the storage key as
the unit of independence with the \emph{lane}. A lane is a
deterministic projection of the touched state onto a domain that is
intuitive to the application, observable to the scheduler, and
finer-grained than the contract address yet coarser than the storage
slot. Lanes make contention \emph{predictable} before execution.

\subsection{Definition}

\begin{definition}[Lane]
\label{def:lane}
A \emph{lane} is the smallest independently schedulable state domain.
Each lane has an identifier
\[
\mathit{lane\_id} =
H\bigl(\mathit{domain},\; \mathit{contract},\; \mathit{account},\;
       \mathit{asset},\; \mathit{market},\; \mathit{nonce\_lane},\;
       \mathit{storage\_prefix}\bigr)
\]
where $H$ is a 256-bit collision-resistant hash. The components are
concrete fields of the touched state; missing components are encoded
as a fixed sentinel.
\end{definition}

\begin{definition}[Lane Hint]
\label{def:lane-hint}
A \emph{lane hint} is a tuple
$\mathit{StmLaneHint} = (\mathit{lane\_id},\, \mathit{access\_kind},\,
\mathit{confidence})$ where
$\mathit{access\_kind} \in \{\textsc{Read},\, \textsc{Write},\,
\textsc{Append},\, \textsc{Reduce},\, \textsc{Unknown}\}$ and
$\mathit{confidence} \in [0,255]$ is the predictor's a priori
estimate that the access will occur.
\end{definition}

A transaction's lane hint vector is computed during admission: by
static ABI declaration when available, by historical observation of
the contract/selector pair otherwise, and by Tier-3 reducer
registration for known commutative operations.

\subsection{Canonical Lane Examples}

\begin{table}[ht]
\centering
\begin{tabular}{@{}ll@{}}
\toprule
\textbf{Lane} & \textbf{Definition} \\
\midrule
sender nonce lane         & $H(\textsf{``nonce''},\, \mathit{account\_id},\, \mathit{nonce\_idx})$ \\
account balance lane      & $H(\textsf{``acct''},\, \mathit{account\_id})$ \\
ERC-20 holder lane        & $H(\textsf{``erc20''},\, \mathit{token},\, \mathit{holder})$ \\
order-book price-level    & $H(\textsf{``ob''},\, \mathit{market\_id},\, \mathit{price\_level})$ \\
margin account lane       & $H(\textsf{``margin''},\, \mathit{account\_id})$ \\
pool reserve lane         & $H(\textsf{``pool''},\, \mathit{pool\_id})$ \\
audit append lane         & $H(\textsf{``audit''},\, \mathit{chain\_id},\, \mathit{epoch})$ \\
fee accumulator lane      & $H(\textsf{``fee''},\, \mathit{token})$ \\
\bottomrule
\end{tabular}
\caption{Canonical lane definitions for the regulated-DEX workload
class. Each lane is a deterministic projection; lane membership is
computable without executing the transaction.}
\label{tab:lane-examples}
\end{table}

\subsection{Lane Clocks}

\begin{definition}[Lane Clock]
\label{def:lane-clock}
A \emph{lane clock} is a monotonic counter
$\mathit{StmLaneClock} = (\mathit{version},\, \mathit{hotness},\,
\mathit{conflict\_count})$ associated with each $\mathit{lane\_id}$.
The \texttt{version} component increments on every committed write
to the lane; the other components are telemetry.
\end{definition}

The clock is a per-lane shard of the TL2 idea. Lane clocks are stored
contiguously in a device-resident hash arena keyed by
$\mathit{lane\_id}$, with collision handling by linear probing under
$\geq 4 \times$ load-factor headroom. There is no global clock
\emph{anywhere} in QuasarSTM.

\subsection{Lane Snapshot}

At execution start a transaction $T$ snapshots the current
\texttt{version} of every lane it touches:
\[
\mathit{LaneSnapshot}(T) = \{(L_i,\, c_i)\}_{i=1}^{m}
\quad \text{where} \quad
c_i = \mathit{LaneClock}[L_i].\mathit{version}.
\]
The snapshot is the precondition for the Tier-1 fast validation
(\S\ref{sec:validation}): if at validation time
$\mathit{LaneClock}[L_i].\mathit{version} = c_i$ for all $i$, the
transaction is committable without further checks.

\subsection{Why Lanes (Not Keys) Are the Right Unit}

\paragraph{Predictability.}
Lane membership is computable from the transaction envelope without
executing the transaction. A scheduler can therefore partition the
ready frontier by lane conflict before any work hits the EVM. Storage
keys are not generally known until execution; they are an outcome,
not a precondition.

\paragraph{Aggregation.}
A single hot lane (\textsc{erc20.totalSupply}, AMM reserve, fee
counter) often subsumes thousands of distinct storage keys whose
semantics commute. Treating these as a single lane and applying a
semantic reducer (\S\ref{sec:validation}) extracts orders of
magnitude more parallelism than per-key MVCC.

\paragraph{Multi-GPU sharding.}
Cross-GPU coherence on storage keys is intractable at scale; the
working set is too fragmented. Lanes give a stable hash domain:
$\mathit{gpu\_id} = H(\mathit{lane\_id}) \bmod \mathit{num\_gpus}$
(\S\ref{sec:gc}) yields a deterministic, well-balanced sharding
plan that survives recompilation.

\paragraph{Telemetry.}
Hot-lane detection
($\mathit{hotness} = \mathit{reads} + 4{\cdot}\mathit{writes} +
8{\cdot}\mathit{conflicts} + 16{\cdot}\mathit{repairs}$)
is naturally per-lane. Per-key hotness aggregates across thousands of
keys before producing a useful signal; per-lane hotness is actionable
within a single round.

\paragraph{Repair priority.}
Cold-lane repairs are background work; hot-lane repairs preempt the
queue. Without a lane abstraction, the scheduler has no signal to
distinguish them.

\begin{remark}[Lane is necessary, not sufficient]
\label{rem:lane-not-sufficient}
Lane prediction does not eliminate MVCC. Misprediction must remain
\emph{safe}: a transaction whose actual access pattern differs from
its hint must still commit a result equivalent to its serial
execution under $\preceq_{\mathrm{can}}$. The Tier-2 key-MVCC
validator (\S\ref{sec:validation}) is the safety net. Lanes are how
QuasarSTM gets fast on the common case; key-MVCC is how it stays
correct on the uncommon case.
\end{remark}


%% ============================================================================
